\subsubsection*{简介}

CDQ 分治是一种思想而不是具体的算法，与动态规划类似。目前这个思想的拓展十分广泛，依原理与写法的不同，大致分为三类：

\begin{itemize}
    \item 解决和点对有关的问题。
    \item 1D 动态规划的优化与转移。
    \item 通过 CDQ 分治，将一些动态问题转化为静态问题。
\end{itemize}

CDQ 分治的思想最早由 IOI2008 金牌得主陈丹琦在高中时整理并总结，它也因此得名。

\hrulefill

\subsubsection*{解决和点对有关的问题}

这类问题多数类似于「给定一个长度为 $n$ 的序列，统计有一些特性的点对 $(i,j)$ 的数量」或「给定一个长度为 $n$ 的序列，找到一对点 $(i,j)$ 使得一些函数的值最大」。

CDQ 分治解决这类问题的算法流程如下：

\begin{enumerate}
    \item 找到这个序列的中点 $mid$；
    \item 将所有点对 $(i,j)$ 划分为 3 类：
    \begin{enumerate}
        \item $1 \le i \le mid,\ 1 \le j \le mid$ 的点对；
        \item $1 \le i \le mid,\ mid+1 \le j \le n$ 的点对；
        \item $mid+1 \le i \le n,\ mid+1 \le j \le n$ 的点对。
    \end{enumerate}
    \item 将 $(1,n)$ 这个序列拆成两个序列 $(1,mid)$ 和 $(mid+1,n)$。此时第一类点对和第三类点对都在这两个序列之中；
    \item 递归地处理这两类点对；
    \item 设法处理第二类点对。
\end{enumerate}

可以看到 CDQ 分治的思想就是不断地把点对通过递归的方式分给左右两个区间。

在实际应用时，我们通常使用一个函数 \texttt{solve(l,r)} 处理 $l \le i \le r,\ l \le j \le r$ 的点对。上述算法流程中的递归部分便是通过 \texttt{solve(l,mid)} 与 \texttt{solve(mid+1,r)} 来实现的。剩下的第二类点对则需要额外设计算法解决。

\paragraph{例题：三维偏序}
给定一个序列，每个点有 $a_i,b_i,c_i$ 三个属性，试求：这个序列里有多少对点对 $(i,j)$ 满足 $a_j \le a_i$ 且 $b_j \le b_i$ 且 $c_j \le c_i$ 且 $j \ne i$。

\textbf{解题思路}：三维偏序是 CDQ 分治的经典问题。题目要求统计序列里点对的个数，那试一下用 CDQ 分治。

首先将序列按 $a$ 排序。

假设已经写好了 \texttt{solve(l,r)}，并且递归搞定了 \texttt{solve(l,mid)} 和 \texttt{solve(mid+1,r)}。现在要做的，就是统计满足 $l \le i \le mid$，$mid+1 \le j \le r$ 的点对 $(i,j)$ 中，有多少点对还满足 $a_i \le a_j$，$b_i \le b_j$，$c_i \le c_j$ 的限制条件。

稍微思考一下就会发现，那个 $a_i \le a_j$ 的限制条件没啥用了：既然 $i$ 比 $mid$ 小，$j$ 比 $mid$ 大，那 $i$ 肯定比 $j$ 要小；已经将序列按 $a$ 排序，就一定有 $a_i \le a_j$。现在还剩下两个限制条件：$b_i \le b_j$ 与 $c_i \le c_j$。根据这个限制条件我们就可以枚举 $j$，求出有多少个满足条件的 $i$。

为了方便枚举，我们把 $(l,mid)$ 和 $(mid+1,r)$ 中的点全部按照 $b$ 的值从小到大排个序。之后我们依次枚举每一个 $j$，把所有 $b_i \le b_j$ 的点 $i$ 全部插入到某种数据结构里（这里我们选择树状数组）。此时只要查询树状数组里有多少个点的 $c$ 值是小于等于 $c_j$ 的，我们就求出了对于这个点 $j$，有多少个 $i$ 可以合法匹配它了。

当我们插入一个 $c$ 值等于 $x$ 的点时，我们就令树状数组的 $x$ 这个位置单点加一，而查询树状数组里有多少个点小于 $x$ 的操作实际上就是在求前缀和。

对于每一个 $j$，我们都需要将所有 $b_i \le b_j$ 的点 $i$ 插入树状数组中。由于所有的 $i$ 和 $j$ 都已事先按照 $b$ 值排好序，这样的话只要以双指针的方式在树状数组里插入点，则对树状数组的插入操作就能从 $O(n^2)$ 次降到 $O(n)$ 次。

通过这样一个算法流程，我们就用 $O(n\log n)$ 的时间处理完了关于第二类点对的信息了。此时算法的时间复杂度是
$$T(n)=T\!\left(\left\lfloor \frac n2 \right\rfloor\right)+T\!\left(\left\lceil \frac n2 \right\rceil\right)+O(n\log n)=O(n\log^2 n)$$

\textbf{实现细节：相等元素要先合并}。若存在三元组完全相同的点，它们之间互相满足偏序关系，但在分治过程中只有「左半区贡献给右半区」这一个方向，相等的点对会漏统计。因此要先把相同的三元组合并成一个点，用 $w$ 记下它的个数，某一点被支配的总数即为 \texttt{res + w - 1}，再把这个数累加到答案数组里（\texttt{ans[d]} 表示被恰好 $d$ 个点支配的点的个数）。

\hrulefill

\subsubsection*{CDQ 分治优化 1D/1D 动态规划的转移}

1D/1D 动态规划指的是一类特定的 DP 问题，该类题目的特征是 DP 数组是一维的，转移是 $O(n)$ 的。如果条件良好的话，有时可以通过 CDQ 分治来把它们的时间复杂度由 $O(n^2)$ 降至 $O(n\log^2 n)$。

例如，给定一个序列，每个元素有两个属性 $a$，$b$。我们希望计算一个 DP 式子的值，它的转移方程如下：

$$dp_i = 1 + \max_{j=1}^{i-1} dp_j\,[a_j < a_i][b_j < b_i]$$

这是一个二维最长上升子序列的 DP 方程，即只有 $j < i,\ a_j < a_i,\ b_j < b_i$ 的点 $j$ 可以更新点 $i$ 的 DP 值。直接转移显然是 $O(n^2)$ 的。

我们发现 $dp_j$ 转移到 $dp_i$ 这种转移关系也是一种点对间的关系，所以我们用类似 CDQ 分治处理点对关系的方式来处理它。假设现在正在处理的区间是 $(l,r)$，算法流程大致如下：

\begin{enumerate}
    \item 如果 $l=r$，说明 $dp_r$ 值的 $\max$ 部分已经被计算好了，直接令 $dp_r \gets dp_r + 1$ 然后返回即可；
    \item 递归使用 \texttt{solve(l,mid)}；
    \item 处理所有 $l \le j \le mid$，$mid+1 \le i \le r$ 的转移关系；
    \item 递归使用 \texttt{solve(mid+1,r)}。
\end{enumerate}

第三步的做法与 CDQ 分治求三维偏序差不多。处理 $l \le j \le mid$，$mid+1 \le i \le r$ 的转移关系的时候，我们会发现已经不用管 $j < i$ 这个限制条件了。因此，我们依然先将所有的点 $i$ 和点 $j$ 按 $a$ 值进行排序处理，然后用双指针的方式将 $j$ 点插入到树状数组里，最后查一下前缀最大值更新一下 $dp_i$ 就可以了。

\paragraph{转移过程的正确性}
该 CDQ 写法和处理点对间关系的 CDQ 写法最大的不同就是处理 $l \le j \le mid$，$mid+1 \le i \le r$ 的点对这一部分。处理点对间关系的 CDQ 写法中，这一部分放到哪里都是可以的。但是，在用 CDQ 分治优化 DP 的时候，这个流程却必须夹在 \texttt{solve(l,mid)}、\texttt{solve(mid+1,r)} 的中间。原因是 DP 的转移是\textbf{有序的}，它必须满足两个条件，否则就是不对的：

\begin{enumerate}
    \item 用来计算 $dp_i$ 的所有 $dp_j$ 值都必须是已经计算完毕的，不能存在「半成品」；
    \item 用来计算 $dp_i$ 的所有 $dp_j$ 值都必须能更新到 $dp_i$，不能存在没有更新到的 $dp_j$ 值。
\end{enumerate}

以 $8$ 这个点为例，它的 DP 值是在 \texttt{solve(1,8)}、\texttt{solve(5,8)}、\texttt{solve(7,8)} 这 3 个函数中更新完成的，而三次用来更新它的点分别是 $(1,4)$、$(5,6)$、$(7,7)$ 这三个不相交的区间。仔细观察就会发现，一个 $i$ 点的 DP 值被更新了 $\log$ 次，而且，更新它的区间刚好是 $(1,i)$ 在线段树上被拆分出来的 $\log$ 个区间。因此，我们的确保证了所有合法的 $j$ 都更新过点 $i$，满足第 2 个条件。

接着分析算法的执行流程：每次 \texttt{solve(l,r)} 结束的时候，$(l,r)$ 区间的 DP 值会被全部计算好。由于我们每一次执行转移函数的时候，\texttt{solve(l,mid)} 已经结束，因此我们每一次执行的转移过程都是合法的，满足第 1 个条件。

如果把 CDQ 分治的递归树看成一颗线段树，那么 CDQ 分治就是这个线段树的\textbf{中序遍历函数}，因此我们相当于按顺序处理了所有的 DP 值，只是转移顺序被拆开了而已，所以算法是正确的。

\hrulefill

\subsubsection*{将动态问题转化为静态问题}

前两种情况使用 CDQ 分治的目的是将序列折半之后递归处理点对间的关系，来获得良好的复杂度。不过在本节中，折半的不是一般的序列，而是\textbf{时间序列}。

它适用于一些「需要支持做 xxx 修改然后做 xxx 询问」的数据结构题。该类题目有两个特点：

\begin{itemize}
    \item 如果把询问离线，所有操作会按照时间自然地排成一个序列。
    \item 每一个修改均与之后的询问操作息息相关。而这样的「修改-询问」关系一共会有 $O(n^2)$ 对。
\end{itemize}

我们可以使用 CDQ 分治对于这个操作序列进行分治，处理修改和询问之间的关系。与处理点对关系的 CDQ 分治类似，假设正在分治的序列是 $(l,r)$，我们先递归地处理 $(l,mid)$ 和 $(mid+1,r)$ 之间的修改-询问关系，再处理所有 $l \le i \le mid$，$mid+1 \le j \le r$ 的修改-询问关系，其中 $i$ 是一个修改，$j$ 是一个询问。

注意，如果各个修改之间是\textbf{独立}的话，我们无需处理 \texttt{solve(l,mid)} 和 \texttt{solve(mid+1,r)} 之间的时序关系（比如普通的加减法问题）。但是如果各个修改之间并不独立（比如说赋值操作），做完这个修改后，序列长什么样可能依赖于之前的序列。此时处理所有跨越 $mid$ 的修改-询问关系的步骤就必须放在 \texttt{solve(l,mid)} 和 \texttt{solve(mid+1,r)} 之间。理由和 CDQ 分治优化 1D/1D 动态规划的原因是一样的：按照中序遍历序进行分治才能保证每一个修改都是严格按照时间顺序执行的。

\paragraph{例题：动态逆序对}
给出 $1 \sim n$ 的一个排列，按照某种顺序依次删除 $m$ 个元素，你的任务是在每次删除一个元素之前统计整个序列的逆序对数。

思路是把删除倒过来看成插入：设第 $i$ 个被删除的元素的出现时间为 $m-i+1$，于是每个元素都有（位置，值，时间）三个属性。一个元素在它被插入时产生的逆序对，就是「位置在它前面且值比它大」以及「位置在它后面且值比它小」的、且时间比它早的元素个数，两者都是三维偏序，用 CDQ 分治统计即可。把答案倒着累加就能得到每次删除前的逆序对数。

\hrulefill

\subsubsection*{习题}

\begin{itemize}
    \item \href{https://www.luogu.com.cn/problem/P3810}{P3810 【模板】三维偏序（陌上花开）}
    \item \href{https://www.luogu.com.cn/problem/P3157}{P3157 [CQOI2011] 动态逆序对}
    \item \href{https://www.luogu.com.cn/problem/P2487}{P2487 [SDOI2011] 拦截导弹}
    \item \href{https://www.luogu.com.cn/problem/P4390}{P4390 [BOI2007] Mokia 摩基亚}
\end{itemize}

\hrulefill

\subsubsection*{参考资料与注释}

\begin{enumerate}
    \item \href{https://oi-wiki.org/misc/cdq-divide/}{OI Wiki：CDQ 分治}
    \item 陈丹琦，\textit{从《Cash》谈一类分治算法的应用}
\end{enumerate}
